\scanpage{143}
\mspage{20}
take $n>1$. First assume that $a\in S$. Then $a=f(x_2/x_1,\ldots,$
$x_n/x_1)$ with $f(Y_2,\ldots,Y_n)\in R[Y_2,\ldots,Y_n]$ so that
$f(x_2/x_1,\ldots,x_n/x_1)=F(x_1,\ldots,x_n)/x_1^q$, where $f(X_1,\ldots,X_n)$
is a form of some degree $q\ge0$ with coefficients in $R$ and
we have $x_1^q a=F(x_1,\ldots,x_n)\in M^q$. Conversely, assume that
$x_1^q a\in M^q$ for some $q\ge0$. Then $x_1^q a=F(x_1,\ldots,x_n)$ where
$F(X_1,\ldots,X_n)$ is a form of degree $q$ with coefficients in $R$
and hence $a=F(x_1,\ldots,x_n)/x_1^q=F(1,x_2/x_1,\ldots,x_n/x_1)\inserted{\in S}$. Now
assume that $a\in x_1S$. Then $a=x_1b$ with $b\in S$ so that
$x_1b\in M^q$ for some $q\ge0$ and hence $x_1^q a=x_1(x_1^q b)\in M^{q+1}$. Con-
versely assume that $x_1^q a\in M^{q+1}$ for some $q\ge0$. Then
$x_1^q a=F(x_1,\ldots,x_n)$ where $F(X_1,\ldots,X_n)$ is a form of degree
$q+1$ with coefficients in $R$ and hence $a/x_1=F(x_1,\ldots,x_n)/x_1^{q+1}=$
$F(1,x_2/x_1,\ldots,x_n/x_1)\in S$ so that $a\in x_1S$.

LEMMA 8. Let $G$ and $H$ be finite cyclic groups and
let $\alpha:G\longrightarrow H$ be an epimorphism. Let $h$ be a
generator of $H$. Then there exists a generator $g$ of
$G$ for which $\alpha(g)=h$.

\uline{Proof}. Let $m$ be the order of $G$ and let $n$ be the order
of $H$. Then $n$ ' divides $m$. Let $r$ be a generator of $G$. Then
for some integer $\deleted{\unsure{m}}\inserted{u}$, $\alpha(r^{\deleted{\unsure{m}}\inserted{\text{\textcircled{$u$}}}})=h$. Let $p_1,\ldots,p_a,q_1,\ldots,q_b$
be the distinct prime divisors of $m$ so labelled that $p_1,\ldots,p_a$
do not divide $u$ and $q_1,\ldots,q_b$ do divide $u$. Let $v=$
$np_1\cdots p_a$, let $w=u+v$ and let $\deleted{g=\unsure{r^m|r^u}}\inserted{g=r^w}$. Since $r$ is a gen-
erator of $G$, $\alpha(r)$ is a generator of $H$. Since $h=\alpha(r^u)=$
$(\alpha(r))^{\deleted{\unsure{m|n}}\inserted{u}}$ is also a generator of $H$ and since $H$ is of order $n$,
\margin{\deleted{\unsure{$r^w$}}}
\margin{$u$}
\figure{\begin{tikzpicture}\draw (0,0)--(.3,-.3)--(.4,.05);\end{tikzpicture}}
\mspage{142}
